% ===== 原文 pp.94--97（tr-07b）=====
% 接 tr07a（p.93 页末）；本片段自 p.94 页首起，末句跨页续入 tr07c 开头。

% [原文 p.94]
至此尚不清楚 $q$-Toda 理论究竟是什么，尤其是：它从根本上说究竟是一个只能放在零曲率曲面上的 2d 理论（这是小弦构造所提示的），还是为了得到有效拉氏量应当把它当作一个 1d 理论 [114]。

\para{椭圆提升。}
把维数再提升一维，$\R^4\times T^2$（扭曲）与 $S^5\times S^1$（超共形指标）上的 6d 配分函数 [646] 与椭圆形变的 $(q,t)$-Toda 有关：见 [272, 306, 647--660]。

\subsection{超共形指标与 2d q-YM}

现在我们转向 4d $\cN=2$ class S 理论在 $M_4=S^3\times S^1$ 上的配分函数。

\para{超对称指标。}
关于超共形 class S 理论见综述 [661]，关于一般 4d $\cN=1$ 理论见 [662]。我们在此不打算多写，而是请读者参看本暑期学校的另一门课程 [663]。与 $q$-YM 的 AGT 关系在 [118] 中也有简要综述。

$S^3\times S^1$ 配分函数对具有无反常 U(1) R 对称性的 4d $\cN=1$ 理论有定义且可计算。除去一个与理论 Casimir 能量（可用 $a$、$c$ 反常表示）有关的因子，该配分函数与超对称指标相符------后者定义为理论在 $S^3\times\R$ 上量子化时的 Witten 指标。一旦用相互对易的旋转、味与 R 对称性的逸度（fugacity）$u_i$（其荷为 $K_i$）加细，指标就写作
\begin{equation}
\mathcal{I}(u)=\Tr\Big[(-1)^F\ee^{-\beta\tilde H}\prod_i u_i^{K_i}\Big]
\tag{9.2}
\end{equation}
其中 $(-1)^F$ 以相反的符号计入玻色态与费米态，$\tilde H=\{Q,Q^\dagger\}$ 对某个超荷 $Q$ 而言，而 $\mathcal I$ 之所以与 $\beta$ 无关，是由于 $\tilde H\neq 0$ 时玻色态与费米态之间的相消。这一简化意味着 $\mathcal I(u)$ 计数（带符号地）超对称代数的短表示。逸度 $u_i$ 在 $S^3\times S^1$ 配分函数中体现为绕 $S^1$ 的、与给定对称性耦合的背景规范场的 holonomy：特别地，两个由旋转与 R 对称性组合而成的逸度 $(p,q)$ 可以理解成 $S^3$ 在 $S^1$ 上的非平凡纤维化。
% [Tutor Note: (9.2) 的核对要点：$\tilde H=\{Q,Q^\dagger\}$ 是超荷的平方（原文 p.94 上文的"超荷平方"），
% 只有在 $\tilde H=0$ 的 BPS 态上 $\ee^{-\beta\tilde H}$ 才不压低，故指标只数短表示、且与 $\beta$ 无关。
% 另注意原文 "it is an renormalization group (RG) flow invariant" 的冠词有笔误（an → a），
% 这是英文语法问题，不影响物理内容。]

指标在形式上不依赖于这些参数之外的任何连续参数：\footnote{在某些语境下，指标的这种形式不变性会失效，由此产生 wall-crossing 现象。}它是一个重整化群（RG）流不变量，例如与规范耦合无关，因此可以在任何弱耦合拉氏量描述中容易地计算。这样一来，$\mathcal I(u)$ 就归结为对任意给定拉氏量描述中由基本场构造的局域规范不变算子作简单的带符号计数（在其他维数中必须把非微扰对象也计入）。作为这样一个极易计算的 RG 流不变量，指标成为窥探 4d $\cN=1$ 规范理论非微扰物理（尤其是其红外对偶）的有力窗口。

% [原文 p.95]
若没有拉氏量描述，计算指标要困难得多，但部分结构仍然保留：若理论 $T$ 是通过对两个理论 $T_1,T_2$ 的公共味对称性 $G$ 作规范化（gauging）而定义的，则指标之间概略地有如下关系
\begin{equation}
\mathcal{I}[T](a_1,a_2)=\int[\dd z]_G\,\mathcal{I}_{\text{vec}}(z)\,\mathcal{I}[T_1](a_1,z)\,\mathcal{I}[T_2](a_2,z),
\tag{9.3}
\end{equation}
其中我们隐去了对 $p,q$ 的依赖，但显式保留了 $T_1$ 与 $T_2$ 的、与 $G$ 对易的味对称性的逸度 $a_1$ 与 $a_2$，它们成为 $T$ 的味对称性。对对称性 $G$ 的逸度 $z$ 的积分用适当的测度 $\mathcal I_{\text{vec}}$ 进行；从局域化的观点看，它就是矢量多重态的单圈行列式。事实上，(9.3) 给出了计算非拉氏量理论指标的一条途径：把它嵌入到一个与拉氏量规范理论对偶的更大理论中，后者的指标是可计算的 [664]。
% [Tutor Note: (9.3) 中记号 $[\dd z]_G$ 是"对规范群 $G$ 的 Cartan 逸度按 Haar 测度积分"的简写，
% $\mathcal I_{\text{vec}}$ 即矢量多重态单圈行列式（局域化视角）。可用 $T_1,T_2$ 均为平凡理论
% （无味对称性、指标为 1）的退化情形核对：此时右边只剩 $\int[\dd z]_G\mathcal I_{\text{vec}}(z)$，
% 应当给出纯矢量多重态理论的指标（即 $G$ 的伴随表示的带符号计数）。]

\para{class S。}
现在我们把讨论专门化到 class S 理论，特别是超共形的那些。由于指标不能依赖规范耦合，它只依赖于黎曼面 $C$ 的拓扑与穿刺点的类型。因此，与标准 AGT 对应相比，2d CFT 一侧应当被一个 TQFT 取代，这一点已在 [665] 中完成。

考虑理论 $\mathrm{T}(\fg,C,D)$。每个穿刺点 $z_i$（$i=1,\dots,n$）联系着一个味对称性，我们打开相应的逸度 $a_i$。对 $C$ 的任意裤子分解，我们都可以把 $\mathrm{T}(\fg,C,D)$ 表示为如下规范化操作的结果：对一族与三穿刺球面相关联的 tinkertoy（孤立的 SCFT）的味对称性作规范化。通过 (9.3)，指标于是归结为 tinkertoy 超共形指标的积分。这恰好模仿了 TQFT 中关联函数的结构：
\begin{equation}
\mathcal{I}\big[T(\fg,C,D)\big](a_i)=\langle\mathcal{O}_{D_1}(a_1)\dots\mathcal{O}_{D_n}(a_n)\rangle_{\text{some TQFT}}
\tag{9.4}
\end{equation}
其中 $\mathcal O_D$ 是依赖于穿刺点类型的适当算子。

与 Liouville CFT 的自举相类似（后者依赖于使用退化顶点算子，它们在规范理论中对应面算子），人们可以用面算子来自举所有 tinkertoy 构造块的指标 [558]。在指标中加入一个面算子，相当于对与任意一个穿刺点相关联的逸度作用一个差分算子 $\Theta$；而由拓扑不变性，是哪一个穿刺点并不重要。把结果在 $\Theta$ 的、以 $\fg$ 的表示 $\lambda$ 标记的本征基中展开，最终得到
\begin{equation}
\mathcal{I}\big[T(\fg,C,D)\big](a_i)=\sum_\lambda (C_\lambda)^{2g-2}\phi^{D_1}_\lambda(a_1)\dots\phi^{D_n}_\lambda(a_n)
\tag{9.5}
\end{equation}
其中 $C_\lambda(p,q,t)$ 是某些结构常数，$\phi^D_\lambda(p,q,t;a)$ 是某些波函数。\footnote{这里我们为 4d $\cN=2$ 理论额外的 R 对称性引入了一个逸度 $t$。}任意穿刺点的波函数可以通过在味逸度中取适当的留数，由满穿刺点的波函数得到 [666]。若穿刺点太少或它们太"小"，这个和可能发散，这或者标志着所给的 class S 理论不存在，或者标志着指标加细得不够------存在额外的味对称性未与任何穿刺点相关联。
% [Tutor Note: (9.5) 中 $(C_\lambda)^{2g-2}$ 的指数可用裤子分解核对：$g=0$、$n=3$ 时
% 该式为 $\sum_\lambda (C_\lambda)^{-2}\phi\phi\phi$，正是 2d TQFT 中把三个波函数收缩成
% 三点关联函数所需的两个逆顶点因子；$g=1$ 时指数为 0，此时环面上的指标不依赖结构常数
% （只有波函数的迹），也与 TQFT 的预期一致。注意：对每个穿刺点都用同一个求和指标 $\lambda$，
% 说明所用基正是 $\Theta$ 对角化后的本征基。]

% [原文 p.96]
这些波函数可以按 $p,q,t$ 逐阶计算，但并不知道其封闭形式。在 Schur 极限 $q=t$ 下，关联函数只是 $q$ 的函数（依赖 $p$ 的项是 $Q$-恰当的），波函数正比于 Schur 多项式，而相应的 TQFT 是 $q$-形变的 2d Yang--Mills 理论 [110]。在更一般的 Macdonald 极限 $p=0$ 下，波函数本质上是 $q,t$ 的 Macdonald 多项式，TQFT 必须通过改变 $q$-YM 理论路径积分中的测度来形变 [667]。Hall--Littlewood 极限 $p=q=0$ 把 Macdonald 多项式变成更简单的 Hall--Littlewood 多项式，后者只依赖 $t$。Coulomb 极限（$t,p\to 0$ 而 $p/t$ 与 $q$ 固定）也很有意思。2d TQFT 描述对一般的 $p,q,t$ 仍相当隐式；它是从 6d 理论 $\cX(\fg)$ 出发在 [668] 中导出的（更早的推导 [669--671] 没有计入瞬子修正）。
% [Tutor Note: p.96 各极限的层次可这样记：$q=t$（Schur）⊂ $p=0$（Macdonald）⊂ 一般 $(p,q,t)$；
% Hall--Littlewood 是 $p=q=0$。核对锚点：$p=0$ 时 Macdonald 多项式退化为 Macdonald 的 $q,t$ 版本，
% 而 $q=t$ 时进一步退化为 Schur 多项式（依赖 $p$ 的项 $Q$-恰当，意味着它们对指标只贡献正合项）。]

这一对应以自然的方式被检验与推广：在 $S^3$ 的极点插入 Wilson--'t Hooft 圈，并相应地在 $q$-YM 中插入圈 [53, 54, 672, 673]；插入一般面算子 [65, 567, 578, 590, 674]；把 $S^3$ 换成透镜空间 $L(p,1)=S^3/\Z_p$ [675--677]；取 $C$ 具有非零面积 [678]；推广到 D 型规范群与非单链（non-simply-laced）的规范群（利用外自同构扭转）[679, 680]。与瞬子模空间 Hilbert 级数的关系在 [216, 226] 中得到探讨。许多 AD 型理论的超共形指标如今也在 Macdonald 极限下已知 [681--685]。这一方向上关键的开放问题似乎是：如何把握完整的参数空间 $(p,q,t)$，而不仅是其 $p=0$ 的 Macdonald 切片。

\subsection{3d/3d 对应}

到目前为止，我们一直在黎曼面上作部分拓扑扭曲，从而约化 $\fg$ 型的 6d $(2,0)$ 理论。若改为约化在 $M_3\times C_3$ 上、沿三维流形 $C_3$ 作扭曲，就得到 $M_3$ 上的 3d $\cN=2$ 规范理论。人们可以加入 6d 理论的余维 2 算子来得到穿刺点的类似物：沿纽结 $K_1\subset C_3$ 的边界条件（在我们的记号中隐去不写）。这定义了一大类 3d $\cN=2$ 规范理论\footnote{这不是标准记号；有时 $T(\fsu(N),C_3,K_1)$ 记作 $T_N[C_3\setminus K_1]$。} $T(\fg,C_3,K_1)$。它们的超对称配分函数与复化 Chern--Simons 理论在 $K_1\subset C_3$ 上的配分函数相符。综述见 [150, 686, 687]；关于该对应及其应用的工作见 [89, 111, 310, 316, 525, 527, 609, 610, 688--727]。

$C_3\setminus K_1$ 的自然构造块是四面体，而 $C_3\setminus K_1$ 的每个三角剖分都给出该 3d $\cN=2$ 理论的一个显式规范理论描述------它是某个阿贝尔 Chern--Simons 理论（以及由质量或其他参数带来的形变）的红外极限。更确切地说，正如 [698, 699] 所指出的，这一描述遗漏了理论 $T(\fg,C_3,K_1)$ 的一部分，而这些额外分支仍在研究中（例如见 [725]）。正如 4d/2d 对应中改变裤子分解等价于 S 对偶，三维流形三角剖分的 Pachner 2-3 移动------把两个相邻四面体换成覆盖流形同一部分的三个四面体------给出 3d $\cN=2$ 的对偶。与 S 对偶不同，这些不是全尺度对偶，而只是红外对偶。
% [Tutor Note: Pachner 2-3 移动（$2\leftrightarrow 3$ 个四面体）与 4d 中改变裤子分解（S 对偶）的类比
% 是本节的结构性要点：两种"改变描述"分别生成 3d N=2 对偶与 4d S 对偶；但前者只在红外成立，
% 因为三角剖分给出的阿贝尔 Chern--Simons 描述本身是 IR 极限，这与正文"not all-scale dualities"一致。]

% [原文 p.97]
\para{对应的陈述。}
3d/3d AGT 对应是在 [111, 525] 中表述的，此前已有若干处理较不一般几何的论文：或者把 $T(\fg,C_3,K_1)$ 进一步在 $S^1$ 上约化 [728]（得到 2d $\cN=(2,2)$ 理论），或者取 $C_3$ 为映射柱或环面（即黎曼面在区间或圆周上的纤维化）[87, 610, 729, 730]。$T(\fg,C_3,K_1)$ 在某些流形 $M_3$ 上的配分函数，等于规范群为 $G_{\C}$（其李代数是 $\fg$ 的复化）的 Chern--Simons 理论在 $C_3\setminus K_1$ 上的配分函数。这反过来提供了纽结与三维流形的不变量。复 Chern--Simons 理论依赖于能级 $(k,\sigma)$：一个量子化的 $k\in\Z$ 和一个连续的 $\sigma\in\R\cup\ii\R$，后者与 $M_3$ 的选取有关。它的作用量形式很直接，
\begin{equation}
S=\frac{k+\ii\sigma}{8\pi}\int_{C_3}\Tr\Big(\mathcal{A}\wedge \dd\mathcal{A}+\frac{2}{3}\mathcal{A}\wedge\mathcal{A}\wedge\mathcal{A}\Big)
+\frac{k-\ii\sigma}{8\pi}\int_{C_3}\Tr\Big(\bar{\mathcal{A}}\wedge \dd\bar{\mathcal{A}}+\frac{2}{3}\bar{\mathcal{A}}\wedge\bar{\mathcal{A}}\wedge\bar{\mathcal{A}}\Big),
\tag{9.6}
\end{equation}
其中 $\mathcal A=A+\ii\Phi$ 是复规范场。
% [Tutor Note: (9.6) 的核对锚点：两个积分项中 $k$ 的系数相加为
% $\frac{k+\ii\sigma}{8\pi}+\frac{k-\ii\sigma}{8\pi}=\frac{k}{4\pi}$，正是标准实 Chern--Simons
% 作用量 $\frac{k}{4\pi}\int\Tr(\mathcal{A}\dd\mathcal{A}+\frac23 \mathcal{A}^3)$ 的系数；
% 而 $\sigma$ 的两个系数 $\frac{\ii\sigma}{8\pi}$ 与 $\frac{-\ii\sigma}{8\pi}$ 相加为 0，即 $\sigma$ 只贡献虚部。
% 若把第二项的 $\bar{\mathcal A}$ 写成 $\mathcal{A}$、或把 $\pm\ii\sigma$ 写成同一符号，实部就不相消，可作为快速判据。]

然而，当规范群非紧时，完整地定义 Chern--Simons 理论是微妙的 [731]；事实上 3d/3d 对应有助于对复规范群 $G_{\C}$ 定义它 [150, 316, 732, 733]。复化 Chern--Simons 理论的少量应用另见 [734, 735]。

压扁球面配分函数（$M_3=S^3_b$）对应能级 $k=1$ 的 Chern--Simons 理论 [111, 310]。超对称指标（$M_3=S^2\times_q S^1$）对应能级 $k=0$ 的 Chern--Simons 理论 [311, 312, 689]。\footnote{我不知道 [311, 312] 之间的差异是否已经解决。}压扁透镜空间 $M_3=L(k,1)_b$ 上的配分函数对应一般 Chern--Simons 能级 $k$ [316, 732]。（这些超对称配分函数以及更多的内容在 [736] 中有综述。）
\begin{equation}
\begin{gathered}
Z_{S^2\times_q S^1}\big[T(\fg,C_3,K_1)\big]=Z_{C_3}\big[G_{\C}\ \text{at levels}\ (0,\sigma)\big],\qquad q=\ee^{2\pi/\sigma}\\
Z_{S^3_b}\big[T(\fg,C_3,K_1)\big]=Z_{C_3}\Big[G_{\C}\ \text{at levels}\ \Big(1,\tfrac{1-b^2}{1+b^2}\Big)\Big]\\
Z_{L(k,1)_b}\big[T(\fg,C_3,K_1)\big]=Z_{C_3}\Big[G_{\C}\ \text{at levels}\ \Big(k,k\tfrac{1-b^2}{1+b^2}\Big)\Big]
\end{gathered}
\tag{9.7}
\end{equation}
% [Tutor Note: (9.7) 三个能级的自洽性核对：$b=1$（正圆 $S^3$）时
% $\frac{1-b^2}{1+b^2}=0$，于是 $S^3_b$ 那一行退化为能级 $(1,0)$ 的普通 Chern--Simons；
% $k=1$ 的透镜空间那一行也同时给出 $(1,0)$（因 $k\frac{1-b^2}{1+b^2}=0$），两行在 $b=1$、$k=1$ 处重合，
% 这正是 $L(1,1)=S^3$ 所要求的。$S^2\times_q S^1$ 的 $q=\ee^{2\pi/\sigma}$ 则把连续能级 $\sigma$
% 与指标中的逸度 $q$ 直接联系起来。]

这三者都可以分解为全纯块（holomorphic blocks）[691, 700, 737]：它们是 Omega 背景 $\R^2\times_q S^1$（或等价地，雪茄 $D^2\times_q S^1$）上的配分函数，在 Chern--Simons 一侧则是波函数。此事的一个半经典版本是：$\R^2\times S^1$ 上超对称真空的集合，与 $C_3\setminus K_1$ 上带有适当边界条件（沿 $K_1$）的平坦 $G_{\C}$ 联络的空间相符。

\para{边界与推广。}
当 $C_3$ 具有 2d 边界（纽结 $K_1$ 可以在边界上终止）时，3d $\cN=2$ 理论生活在与 $\partial C_3$ 相关联的 4d $\cN=2$ class S 理论的边界上（并与它耦合）[43, 89, 111, 673]。特别地，当 $C_3$ 是一个配边（cobordism），即 $\partial C_3$ 由两个不连通分支组成时，更自然的想法是把 3d $\cN=2$ 理论看作相应的两个 4d $\cN=2$ class S 理论之间的畴壁，而这两个理论只通过它们共同的 3d 边界耦合。于是该构造对于粘合是函子性的。关于这一设置的一个特别简单的例子在 [43] 中有描述：取 $C_3=\R\times C$，其中 $C$ 的复结构沿 $\R$ 方向变化；那么在规范理论一侧，我们有 4d $\cN=2$ 理论
